\subsection{Memoria acumulada y resolventes del transporte}
\label{subsec:memoria-resolvente}

El retorno del calendario no elimina los incrementos producidos durante
el recorrido. Para expresar esta diferencia operatoriamente, se construye
primero un registro acumulativo de eventos de la trayectoria y después su
serie generadora. La reducción de variables de memoria se realizará sobre
ese transporte, conservando también la fuente y el lector de salida.

\subsubsection{Eventos orientados y actualización del registro}

Sea \((d_t)\) la secuencia de códigos enteros de dirección conservada
por una trayectoria del TPK. El código \(d_t\) y la dirección cardinal
del cursor son coordenadas distintas; no se supone una identificación
entre sus alfabetos. Numeramos desde cero las ventanas
\[
 I_m=\{9m+1,\ldots,9m+9\},\qquad m\geq0.
\]
En el primer ciclo, \(I_m\) es la ventana \(E_{m+1}\) del registro
dodecafásico. Su orientación es
\(\sigma=(1,1,1,-1,-1,-1,1,1,1,-1,-1,-1)\).
Para una historia prolongada, \(\sigma_m\in\{-1,1\}\) es el signo
conservado en la ventana correspondiente. El evento firmado es
\begin{equation}
 a_m=\sigma_m\sum_{t\in I_m}
       \mathbf1_{\{2,4,6,8\}}(d_t),\qquad |a_m|\leq9.
 \label{mem:evento}
\end{equation}
Cada sumando se incorpora cuando ocurre su transición. Esta definición
cuenta los eventos de la traza; no escoge direcciones mediante un valor
que deba obtenerse después.

Sean \(\mathbf e_0,\ldots,\mathbf e_{11}\) la base de \(\ZZ^{12}\)
y \(S\mathbf e_j=\mathbf e_{j+1\bmod12}\). Escribimos
\(R=S^{-1}\). El registro acumulado en coordenadas fijas satisface
\[
z_{m+1}=z_m+a_m\mathbf e_{m\bmod12}.
\]
Un prefijo sin historia acumulada empieza en \(z_0=0\); en otro
origen, \(z_0\) conserva el registro previo.
El marco que acompaña al calendario es \(y_m=S^{-m}z_m\).

\begin{proposition}[Actualización y retorno con memoria]
\label{mem:prop-actualizacion}
El registro en el marco móvil cumple
\begin{equation}
 y_{m+1}=Ry_m+\eta_m,\qquad
 \eta_m=a_mR\mathbf e_0.
 \label{mem:actualizacion}
\end{equation}
Para todo \(n\geq1\),
\begin{equation}
 y_n=R^ny_0+\sum_{j=0}^{n-1}R^{n-1-j}\eta_j.
 \label{mem:iteracion}
\end{equation}
En particular,
\begin{equation}
 y_{12}=y_0+\sum_{j=0}^{11}a_j\mathbf e_j.
 \label{mem:retorno12}
\end{equation}
\end{proposition}
\begin{proof}
Multiplicar la actualización de \(z_m\) por \(S^{-(m+1)}\) da
\(y_{m+1}=S^{-1}y_m+a_mS^{-1}\mathbf e_0\), porque
\(S^{-m}\mathbf e_{m\bmod12}=\mathbf e_0\).
La fórmula \eqref{mem:iteracion} se obtiene por inducción.
Para \(n=12\), \(R^{12}=I\) y
\(R^{11-j}\eta_j=a_jR^{12-j}\mathbf e_0=a_j\mathbf e_j\),
lo que prueba \eqref{mem:retorno12}.
\end{proof}

Así, doce ventanas restauran el marco y conservan el vector de eventos.
El recuento no permite recuperar por sí solo el orden de todos los pasos
dentro de cada ventana; ese orden permanece en la trayectoria registrada.

Para precisar el balance, definimos
\[
 D_3=I-S^3,\quad D_4=I-S^4,\quad
 B=D_3^{\mathsf T}D_3+D_4^{\mathsf T}D_4,
\]
\[
 \mathcal M(y)=\tfrac12y^{\mathsf T}By,
 \qquad q(y)=\sum_{j=0}^{11}y_j.
\]
La primera es una forma cuadrática del registro; la segunda registra
la suma de sus coordenadas. No se identifican aquí con energía ni carga físicas.

\begin{proposition}[Balance producido por el evento]
\label{mem:prop-balance}
La actualización \eqref{mem:actualizacion} satisface
\begin{equation}
 \begin{aligned}
 \mathcal M(y_{m+1})-\mathcal M(y_m)
   &=a_m(By_m)_0+2a_m^2,\\
 q(y_{m+1})-q(y_m)&=a_m.
 \end{aligned}
 \label{mem:balance}
\end{equation}
\end{proposition}
\begin{proof}
Al desarrollar los productos,
\(B=4I-S^3-S^{-3}-S^4-S^{-4}\); por tanto
\(R^{\mathsf T}BR=B\) y \(B_{00}=4\).
Como \(y_{m+1}=R(y_m+a_m\mathbf e_0)\), la diferencia cuadrática
es \(a_m\mathbf e_0^{\mathsf T}By_m+\tfrac12a_m^2B_{00}\).
La suma de coordenadas es invariante bajo la permutación \(R\),
lo que da la segunda igualdad.
\end{proof}

\Needspace{20\baselineskip}
\subsubsection{Serie generadora y control uniforme de la memoria}

La serie conserva todos los incrementos, no únicamente su suma al final
de un ciclo. Con una variable formal \(u\), sean
\[
 Y(u)=\sum_{m\geq0}y_mu^m,\qquad
 H_\eta(u)=\sum_{m\geq0}\eta_mu^m.
\]

\begin{proposition}[Resolvente y cota de cola]
\label{mem:prop-serie}
Se cumple la identidad formal
\begin{equation}
 Y(u)=(I-uR)^{-1}\bigl(y_0+uH_\eta(u)\bigr).
 \label{mem:serie}
\end{equation}
Ambas series convergen para \(|u|<1\). Si \(\ell:\RR^{12}\to\RR\)
es lineal, \(L=\|\ell\|_{1\to\RR}\), \(M_0=\|y_0\|_1\) y
\(0<r<1\), entonces, para todo entero \(N\geq0\) y \(|u|\leq r\),
\begin{equation}
 \left|\sum_{m>N}\ell(y_m)u^m\right|
 \leq Lr^{N+1}\left(
 \frac{M_0+9(N+1)}{1-r}+\frac{9r}{(1-r)^2}\right).
 \label{mem:cota}
\end{equation}
La serie de derivadas converge uniformemente en cada disco
\(|u|\leq r<1\).
\end{proposition}
\begin{proof}
Multiplicar \eqref{mem:actualizacion} por \(u^{m+1}\) y sumar da
\(Y-y_0=uRY+uH_\eta\). El término constante de \(I-uR\) es
invertible, de modo que su inversa formal es \(\sum_{k\geq0}u^kR^k\).
Además, \(R\) preserva la norma \(\ell^1\) y
\(\|\eta_m\|_1\leq9\), de donde
\(\|y_m\|_1\leq M_0+9m\).
Para la cola se suman las mayorantes
\(L(M_0+9m)r^m\), usando
\[
 \sum_{m>N}r^m=\frac{r^{N+1}}{1-r},\qquad
 \sum_{m>N}mr^m=r^{N+1}
 \left(\frac{N+1}{1-r}+\frac r{(1-r)^2}\right).
\]
Por último, \(\sum_{m\geq1}Lm(M_0+9m)r^{m-1}<\infty\)
mayora la serie derivada y prueba su convergencia uniforme.
\end{proof}

La variable \(u\) cuenta ventanas de nueve transiciones. Su identificación
con la variable de otro lector requiere conservar ese reloj y el mapa
que relaciona ambas lecturas.

\subsubsection{Eliminación de memoria: operador, fuente y lector}

En un dominio donde la actualización está fijada como \(U:X\to X\),
el espacio vectorial libre \(V=\QQ^{(X)}\) sobre los estados admite el operador
\(T\delta_x=\delta_{U(x)}\). El reloj se incluye en \(x\) si la
transición depende de él. Esta extensión lineal conserva estados
distintos aunque compartan una fase observable.

Consideremos una descomposición \(V=V_0\oplus V_1\), donde \(V_1\)
es el sector auxiliar que se desea eliminar, y escribamos
\[
 T=\begin{pmatrix}A&B_1\\C&D\end{pmatrix},\qquad
 v=\binom{v_0}{v_1},\qquad \ell=(\ell_0,\ell_1).
\]
Aquí \(A:V_0\to V_0\), \(B_1:V_1\to V_0\),
\(C:V_0\to V_1\) y \(D:V_1\to V_1\).
Las identidades siguientes son formales; también valen donde los
resolventes convergen.

\begin{proposition}[Reducción de Schur con fuente y lector]
\label{mem:prop-schur}
Sean
\[
 D_u=(I-uD)^{-1},\qquad
 \mathscr F(u)=I-uA-u^2B_1D_uC.
\]
La componente retenida de \((I-uT)^{-1}v\) es
\begin{equation}
 w_0=\mathscr F(u)^{-1}\bigl(v_0+uB_1D_uv_1\bigr),
 \qquad w_1=D_u(v_1+uCw_0).
 \label{mem:schur-fuente}
\end{equation}
Su lectura completa satisface
\begin{equation}
 \begin{split}
 \ell(I-uT)^{-1}v
 ={}&(\ell_0+u\ell_1D_uC)
       \mathscr F(u)^{-1}(v_0+uB_1D_uv_1)\\
    &+\ell_1D_uv_1.
 \end{split}
 \label{mem:schur-lector}
\end{equation}
\end{proposition}
\begin{proof}
Las dos ecuaciones de \((I-uT)w=v\) son
\(w_0=v_0+uAw_0+uB_1w_1\) y
\(w_1=v_1+uCw_0+uDw_1\).
Resolver la segunda y sustituirla en la primera da
\eqref{mem:schur-fuente}; aplicar \(\ell_0\) y \(\ell_1\)
a ambas componentes da \eqref{mem:schur-lector}.
Todas las inversas formales existen porque sus términos constantes
son identidades.
\end{proof}

El término \(u^2B_1D_uC=\sum_{k\geq0}u^{k+2}B_1D^kC\)
describe una salida hacia la memoria, \(k\) pasos dentro de ella y
un regreso. En cambio, \(u\ell_1D_uC\) lee la memoria antes de ese
regreso. Por ello, reducir sólo el operador no conserva necesariamente
la observación: también deben transformarse fuente y lector.

La diferencia es visible en el ciclo ya construido. Si \(P\) retiene
\(\mathbf e_0\), entonces \(PRP=0\), pero
\begin{equation}
 \left\langle\mathbf e_0,(I-uR)^{-1}\mathbf e_0\right\rangle
 =\frac1{1-u^{12}},\qquad
 \left\langle\mathbf e_{11},(I-uR)^{-1}\mathbf e_0\right\rangle
 =\frac u{1-u^{12}}.
 \label{mem:ejemplo-ciclo}
\end{equation}
En efecto, \(R^n\mathbf e_0=\mathbf e_{-n\bmod12}\): las
dos series recogen, respectivamente, \(n\equiv0\) y
\(n\equiv1\pmod{12}\). El resolvente de la compresión es \(I\),
mientras la compresión del resolvente conserva los retornos omitidos.
Este ejemplo corresponde al registro de doce ventanas, no a la totalidad
del estado del TPK.

\subsubsection{Composición de segmentos y coherencia de la reducción}

Sean tres transportes afines composables
\(F_i(y)=R_i y+b_i\), con dominios y codominios compatibles.
La memoria producida por el recorrido completo es
\begin{equation}
 F_3F_2F_1(y)
 =R_3R_2R_1y+b_3+R_3b_2+R_3R_2b_1.
 \label{mem:cociclo-tres}
\end{equation}
La prueba es la sustitución sucesiva de las tres fórmulas afines.
Agrupar primero los dos primeros segmentos o los dos últimos produce
la misma expresión; los factores transportan cada incremento desde
el punto donde se originó. En \eqref{mem:actualizacion},
\(R_i=R\) y \(b_i\) es el evento de su ventana.

La eliminación de varias capas conserva una coherencia semejante.
Para hacerla explícita, descomponemos ahora
\(V=V_0\oplus V_1\oplus V_2\) y escribimos
\[
 T=\begin{pmatrix}
 A&B_1&B_2\\C_1&D_{11}&D_{12}\\C_2&D_{21}&D_{22}
 \end{pmatrix},\qquad Q_2=(I-uD_{22})^{-1}.
\]
Al eliminar la tercera componente quedan los cuatro bloques
\begin{equation}
 \begin{aligned}
 A'&=A+uB_2Q_2C_2,&
 B_1'&=B_1+uB_2Q_2D_{21},\\
 C_1'&=C_1+uD_{12}Q_2C_2,&
 D_{11}'&=D_{11}+uD_{12}Q_2D_{21}.
 \end{aligned}
 \label{mem:schur-intermedio}
\end{equation}

\begin{proposition}[Transitividad de la eliminación]
Si \(D_*=(D_{ij})_{i,j=1}^2\), el complemento obtenido en dos etapas
coincide con el obtenido al eliminar juntas ambas memorias:
\begin{equation}
 \begin{split}
 I-uA'-u^2B_1'(I-uD_{11}')^{-1}C_1'
 ={}&I-uA\\
 &-u^2(B_1\ B_2)(I-uD_*)^{-1}\binom{C_1}{C_2}.
 \end{split}
 \label{mem:schur-transitivo}
\end{equation}
\end{proposition}
\begin{proof}
Resolver la tercera ecuación de \((I-uT)w=v\) y sustituirla en las
dos primeras produce \eqref{mem:schur-intermedio}. Eliminar después la
segunda componente da el miembro izquierdo de
\eqref{mem:schur-transitivo}. Resolver a la vez las dos ecuaciones de
memoria da el derecho, por la proposición~\ref{mem:prop-schur}.
La solución formal es única porque \(I-uT\) tiene término constante
\(I\). En ambos procedimientos se transforman también la fuente y
el lector según \eqref{mem:schur-fuente}--\eqref{mem:schur-lector}.
\end{proof}

El balance \eqref{mem:balance} describe los incrementos del registro.
Las identidades anteriores permiten transportarlos al reutilizar la memoria
en la clausura dodecafásica; no determinan por sí solas coeficientes
analíticos adicionales.
